\section{Introduction}
Text-to-brick generation connects natural-language descriptions with the spatial arrangement of discrete building pieces. A concise request such as \emph{a simple chair} identifies an object category while leaving its dimensions, proportions, and part layout largely unspecified. A generator must complete these geometric details and express the resulting shape through brick placements. Addressing this challenge involves both text--geometry supervision spanning different amounts of descriptive information and explicit constraints governing how discrete pieces connect.

BrickGPT~\cite{brickgpt} learns autoregressive brick-placement sequences and incorporates validity checks with physics-aware rollback during inference. BrickNet~\cite{bricknet} adopts a graph-based representation supporting a broader range of parts and connection types. Alongside these advances, generation from concise descriptions warrants explicit examination. In our five-level benchmark (Table~\ref{tab:level-results}), BrickNet obtains lower mean VQAScore, PE, and SigLIP2 than BrickForma at the two shortest description levels, L1 and L2. This observed gap motivates a closer examination of how generators translate limited linguistic information into object geometry. Level-wise evaluation makes performance on concise requests visible alongside performance on detailed specifications.

We present BrickForma, placing programmatic geometric data construction at the foundation of text-conditioned brick generation. With large-language-model assistance, we develop parametric programs that vary object proportions, part dimensions, openings, and support arrangements to synthesize voxel shapes in batches. The source collection contains 80,000 models across 21 categories on a $48^3$ grid. Each model is paired with five levels of bilingual description that progressively introduce structural, geometric, proportional, and appearance information. These approximately 400,000 caption records link concise requests and detailed specifications to shared geometry, providing multilevel supervision for text-conditioned spatial layout prediction.

Converting voxel shapes into brick targets explicitly addresses the distinction between spatial adjacency and stud connectivity. A connected voxel volume can be tiled by pieces that touch side by side without connecting through studs. Our conversion pipeline combines same-color tiling, connection analysis, and local bridging and retiling to construct geometrically checked target layouts, with changes to source voxels recorded in derived models. The integrated dataset release includes 79,301 models passing its dataset-level geometric checks. We fine-tune Qwen3-1.7B~\cite{qwen3} with low-rank adaptation~\cite{lora} on caption--placement pairs to predict brick dimensions and spatial positions.

Geometric reconstruction complements the language model's placement predictions. From the supported, in-bounds portions of predicted placements, the pipeline forms a voxel union and uses mixed-integer linear programming to seek an exact-cover layout forming a single stud-connected component. Independent coverage and connectivity checks verify the reconstructed result. This design decouples text-conditioned shape prediction in the language model from explicit structural enforcement in the reconstruction stage.

We evaluate BrickForma against BrickGPT and BrickNet on 114 shared prompts spanning five description levels. BrickForma achieves the highest mean VQAScore, PE, and SigLIP2 at every level, including the shortest descriptions, while matching BrickGPT's mean invalid-placement ratio of 0.00\%. Reporting all five levels connects the framework's multilevel data design to an explicit assessment of generation across different amounts of descriptive detail.

Our contributions are:
\begin{enumerate}
\item An LLM-assisted parametric data-construction pipeline producing 80,000 source voxel models with five levels of associated description, linking varying amounts of linguistic detail to shared object geometry.
\item A text-to-brick framework integrating voxel-to-brick target construction, Qwen3-1.7B fine-tuning, and MILP-based exact-cover reconstruction with independent coverage and stud-connectivity checks.
\item A five-level evaluation against BrickGPT and BrickNet that explicitly examines concise descriptions alongside detailed specifications through text--image alignment and invalid-placement measurements.
\end{enumerate}
